\chapter{Detalle de la iteración}

\section{Historias de usuario}
La siguiente tabla contiene las actividades mayores a realizar 
durante la iteración.

\subsection{Descripción de las columnas}
\begin{itemize}
  \item \textbf{ID}: El identificador del Hito, seguido por el identificador de la Historia.
  \item \textbf{Historia}: Título del Hito, seguido por el título de la Historia.
  \item \textbf{Estado}: Estado actual de la Historia.
  \item \textbf{Término}: Fecha límite para completar la Historia.
\end{itemize}

\input{generated/stories.tex}

\pagebreak

\section{Sub-tareas}
La siguiente tabla contiene las actividades menores a realizar 
durante la iteración.

\subsection{Descripción de las columnas}
\begin{itemize}
  \item \textbf{ID}: El identificador de la Historia padre, seguido por el identificador de la Sub-tarea.
  \item \textbf{Sub-tarea}: Título de la Historia padre, seguido por el título de la Sub-tarea.
  \item \textbf{Estado}: Estado actual de la Sub-tarea.
  \item \textbf{Término}: Fecha límite para completar la Sub-tarea.
\end{itemize}

\input{generated/subtasks.tex}

\pagebreak

\section{Tareas pendientes}
Las siguientes tareas \textbf{no} se incluirán en la iteración actual, pero se 
espera que sean completadas en iteraciones futuras.

\subsection{Descripción de las columnas}
\begin{itemize}
  \item \textbf{ID}: El identificador de la Historia padre, seguido por el identificador de la Sub-tarea.
  \item \textbf{Sub-tarea}: Título de la Historia padre, seguido por el título de la Sub-tarea.
  \item \textbf{Estado}: Estado actual de la Sub-tarea.
\end{itemize}

\input{generated/backlog.tex}